\section{Experimental setup}
\label{sec:experiments}

\subsection{Data}

We collect $(s_t, a_t, s_{t+1}, w_{t+1})$ transitions for each game by
running the C++ PuzzleScript engine under bounded A* search from each
authored level. Each transition is cached on disk (bit-packed for
roughly 8$\times$ compression) and re-used across runs. For
multi-game training we draw batches with one share per game per step
(\emph{balanced sampling}), which compensates for per-game dataset-size
imbalance.

\paragraph{Synthetic levels.}
For ablations and held-out evaluations we additionally generate
synthetic levels by sampling tile-pattern distributions from each
game's authored levels and validating each candidate by running the
engine to confirm the level is solvable and exposes a non-trivial
state space (more details in the appendix). The validated synthetic
recipe (per-game sizes, fallback dynamics, multi-grid widths) is what
allowed our model to match the authored-data baseline on the
\textsc{nekopuzzle} game without overfitting to grid-size artifacts.

\subsection{Game suites}

We evaluate on three suites of increasing breadth:

\begin{itemize}
  \item \textbf{Single-game} (\textsc{microban}, 10 authored levels):
    classic chained Sokoban pushing on small grids. Used for
    architecture ablations where dynamics complexity is fully
    controlled.
  \item \textbf{Single-game / synthetic-multi-grid}
    (\textsc{sokoban\_basic}): a controlled stress test in which the
    model is trained on synthetic levels of multiple grid sizes
    $\{5{\times}5, 6{\times}6, 7{\times}7, 8{\times}8\}$ and evaluated
    on the same. The diversity prevents per-state memorization and
    forces the model to learn the underlying rule.
  \item \textbf{Multi-game} (\textsc{scaling\_14}): a fourteen-game
    benchmark including chained pushing, gravity, light beams,
    multi-character control, and several \texttt{again}-driven
    simulations. Per-game grid sizes range from $4{\times}4$ to
    $32{\times}24$.
\end{itemize}

\subsection{Architectures and hyper-parameters}

All models share the same input embedding (Dense over
$(s_t \,\|\, \mathrm{onehot}(a_t))$) and readout heads. The slot
encoder is identical across architectures: $K{=}16$ slots, $d{=}64$
slot dim, two self-attention layers over tokens, four heads.

The NCA uses $T{=}4$ NCA steps with $n_r{=}T$ (single shared layer
applied $T$ times), $d_h{=}256$, full pool features, input-skip on,
optional padding mask for multi-grid. The CNN baseline uses
$n_b{=}4$ residual blocks; the U-Net uses 2 levels of down/up
sampling with bottleneck FiLM; the ViT uses $n_v{=}4$ layers, four
heads. We report parameter count alongside every metric so the
comparison is honest.

Training: Adam, learning rate $3 \times 10^{-4}$ with cosine decay
to $10^{-7}$, gradient clipping at global norm $0.5$, batch size 32,
\emph{change-loss-weight} $5.0$ (per-cell BCE upweighted on cells
that change between $t$ and $t+1$), and patience-based early stop on
the change-accuracy metric. Each run is replicated with three random
seeds.

\subsection{Evaluation}

We report two complementary metrics:
\begin{itemize}
  \item \textbf{One-step change-error}: the per-cell BCE-thresholded
    error rate on the cells that change between $s_t$ and $s_{t+1}$,
    measured under teacher forcing on the held-out portion of each
    cache.
  \item \textbf{Autoregressive rollout error}: the same per-cell error
    averaged over a 50-step rollout from each level's start state,
    under three policies: random actions, BFS-optimal, and A*-optimal.
\end{itemize}
Per finding F3 in our architecture report~(\S\ref{sec:results}), train
loss decouples from autoregressive rollout error in the
over-capacity regime; reporting both is essential.
